\section{从结构开始}\label{sec:structure}
正文只描述内容的层级与含义。章、节、图、表、公式的序号由计数器管理，
引用目标由标签标识。不要用手打空格对齐正文，也不要把编号写死在句子中。
例如数学部分见\cref{sec:math}，坐标图见\cref{fig:coordinate-pair}。

\begin{itemize}
  \item \texttt{main.tex}：文档类、宏包、字体、页眉与章节装配。
  \item \texttt{sections/}：按学习主题组织的正文源文件。
  \item \texttt{figures/}：坐标图的 TikZ 源文件及编译生成的 PDF。
  \item \texttt{references.bib}：真实文献的元数据，引用键不依赖显示顺序。
\end{itemize}

\textbf{第一项练习：}在本节前增加一个小节，重新构建后检查目录、正文引用与页码。
一次构建应让依赖稳定；PDF 能打开，只能说明产物存在，不能替代检查警告和版面。
学习排版底层模型可以参阅 Knuth 的著作\cite{knuth1984texbook}。
